\documentclass{beamer}
\usetheme{Madrid}
\usecolortheme{default}
\usepackage[utf8]{inputenc}
\usepackage{amsmath}
\usepackage{amsfonts}
\usepackage{amssymb}
\usepackage{graphicx}
\usepackage{array}
\usepackage{colortbl}
\usepackage{multirow}
\usepackage{tikz}
\usepackage{xcolor}

\title{AML/CFT Certification: Difficult Topics \& Answer Strategies}
\subtitle{Comprehensive Exam Preparation Guide}
\author{Compliance Professional Training}
\date{Updated: June 2026}

\definecolor{correct}{RGB}{0,128,0}
\definecolor{wrong}{RGB}{220,20,60}
\definecolor{highlight}{RGB}{255,215,0}
\definecolor{info}{RGB}{0,102,204}
\definecolor{warning}{RGB}{255,140,0}

\setbeamercolor{block title}{bg=info!20,fg=info}
\setbeamercolor{block body}{bg=info!5}
\setbeamercolor{alerted text}{fg=wrong}

\begin{document}

% ================================
% SLIDE 1: TITLE
% ================================
\begin{frame}
\titlepage
\end{frame}

% ================================
% SLIDE 2: TABLE OF CONTENTS
% ================================
\begin{frame}{Table of Contents}
\tableofcontents
\end{frame}

% ================================
% SLIDE 3: CORE THEMES
% ================================
\section{Core Themes}
\begin{frame}{Core Themes Across All Exam Questions}
\begin{block}{5 Key Concepts You Must Master}
\begin{enumerate}
    \item \textbf{Risk-Based Approach} - Action proportional to risk level
    \item \textbf{Documentation \& Audit Trail} - Record every step
    \item \textbf{Technology \& Automation} - Systems require tuning
    \item \textbf{International Standards} - FATF = Global baseline
    \item \textbf{Beneficial Ownership} - Know ultimate controller
\end{enumerate}
\end{block}
\begin{alertblock}{Exam Mindset}
\textbf{Ask yourself:} \textit{"What would a regulator expect?"} \\
\textbf{Ask yourself:} \textit{"What would an auditor check?"}
\end{alertblock}
\end{frame}

% ================================
% SLIDE 4: TRANSACTION MONITORING - OVERVIEW
% ================================
\section{Transaction Monitoring}
\begin{frame}{Transaction Monitoring: The Process}
\begin{block}{Standard Alert Review Workflow}
\begin{enumerate}
    \item System triggers alert on transaction
    \item Analyst reviews transaction patterns
    \item Analyst gathers information from sources
    \item Documents findings throughout investigation
    \item Escalates if suspicious activity confirmed
\end{enumerate}
\end{block}
\begin{alertblock}{Golden Rule}
\textbf{Never rely solely on automated systems -} human judgment is essential!
\end{alertblock}
\end{frame}

% ================================
% SLIDE 5: TRANSACTION MONITORING - QUESTION 1
% ================================
\begin{frame}{Transaction Monitoring: Information Sources}
\begin{block}{The Question}
\small{While reviewing an alert, the analyst notes wire transfers to an unfamiliar entity in a sanctioned jurisdiction. The customer profile does not list dealings in that region. Which information source would most help the analyst understand the activity?}
\end{block}

\vspace{0.3cm}
\begin{columns}
\column{0.5\textwidth}
\textcolor{wrong}{\textbf{Wrong Answer:}}
\begin{itemize}
    \item Regulatory advisories about common evasion techniques in that jurisdiction
\end{itemize}

\column{0.5\textwidth}
\textcolor{correct}{\textbf{right Correct Answer:}}
\begin{itemize}
    \item Open-source research into the entity
\end{itemize}
\end{columns}

\vspace{0.3cm}
\begin{block}{Why?}
Specific entity research \textgreater{} General guidance. \\
\textbf{Understand WHO you're dealing with first.}
\end{block}
\end{frame}

% ================================
% SLIDE 6: TRANSACTION MONITORING - QUESTION 2
% ================================
\begin{frame}{Transaction Monitoring: Documentation Best Practices}
\begin{block}{The Question}
\small{Which of the following best practices support effective research documentation in suspicious activity investigations? (Select Two.)}
\end{block}

\begin{columns}
\column{0.5\textwidth}
\textcolor{wrong}{\textbf{ Wrong:}}
\begin{itemize}
    \item Documenting findings at the conclusion of an investigation
\end{itemize}

\column{0.5\textwidth}
\textcolor{correct}{\textbf{Correct:}}
\begin{itemize}
    \item Using structured templates for investigative notes
    \item Adapting tools based on jurisdictional requirements
\end{itemize}
\end{columns}

\vspace{0.3cm}
\begin{block}{Key Takeaway}
Document \textbf{throughout}, not just at end. \\
Use \textbf{structured templates} for consistency.
\end{block}
\end{frame}

% ================================
% SLIDE 7: TRANSACTION MONITORING - QUESTION 3
% ================================
\begin{frame}{Transaction Monitoring: Frontline Escalation}
\begin{block}{The Question}
\small{A frontline staff receives a large, unusually timed cash deposit from a long-standing customer at a small bank branch. There is no immediate system alert. What is the most appropriate action?}
\end{block}

\begin{columns}
\column{0.5\textwidth}
\textcolor{wrong}{\textbf{ Common Trap:}}
\begin{itemize}
    \item Ignore because no system alert
    \item Wait for system to catch it
\end{itemize}

\column{0.5\textwidth}
\textcolor{correct}{\textbf{Correct Answer:}}
\begin{itemize}
    \item Manually escalate a report to compliance
\end{itemize}
\end{columns}

\begin{block}{Why This Matters}
\textbf{Human judgment > Automated systems} \\
Unusual patterns below threshold still need review!
\end{block}
\end{frame}

% ================================
% SLIDE 8: PRIVATE BANKING RISKS
% ================================
\section{Private Banking Risks}
\begin{frame}{Private Banking: High-Risk Features}
\begin{block}{Question: High-Risk Features of Private Banking (Select Three)}
\end{block}

\begin{columns}
\column{0.5\textwidth}
\textcolor{wrong}{\textbf{ Wrong Answers:}}
\begin{itemize}
    \item Use of free-trade zones
\end{itemize}

\column{0.5\textwidth}
\textcolor{correct}{\textbf{Correct Answers:}}
\begin{itemize}
    \item Minimum assets requirement
    \item Offering numbered accounts
    \item Dedicated relationship service with transaction discretion
\end{itemize}
\end{columns}

\vspace{0.3cm}
\begin{block}{Why These Are High-Risk}
\begin{itemize}
    \item \textbf{Minimum assets} - Attracts high-net-worth, complex structures
    \item \textbf{Numbered accounts} - Reduced transparency
    \item \textbf{Relationship discretion} - Potential to bypass controls
\end{itemize}
\end{block}
\end{frame}

% ================================
% SLIDE 9: PRIVATE BANKING - THEORY
% ================================
\begin{frame}{Private Banking Risk Theory}
\begin{block}{Why Private Banking is Vulnerable}
\begin{itemize}
    \item \textbf{Personalized relationship management} creates:
    \begin{itemize}
        \item Pressure to maintain client relationships
        \item Potential for "turning a blind eye"
        \item Discretion that may hide activity
    \end{itemize}
    \item \textbf{Offshore jurisdictions} complicate oversight
    \item \textbf{Complex structures} obscure beneficial ownership
\end{itemize}
\end{block}

\begin{alertblock}{Exam Tip}
Questions about private banking = think about \textbf{discretion, opacity, and relationship pressure}.
\end{alertblock}
\end{frame}

% ================================
% SLIDE 10: CONCENTRATION ACCOUNTS
% ================================
\begin{frame}{Concentration Accounts: Expected Uses}
\begin{block}{Question: Typical and Expected Uses of Concentration Accounts (Select Two)}
\end{block}

\begin{columns}
\column{0.5\textwidth}
\textcolor{wrong}{\textbf{ Wrong:}}
\begin{itemize}
    \item Aggregating transactions from multiple unrelated customers
\end{itemize}

\column{0.5\textwidth}
\textcolor{correct}{\textbf{Correct:}}
\begin{itemize}
    \item Streamlining inter-branch cash movements within the bank
    \item Facilitating back-office functions like suspense clearing
\end{itemize}
\end{columns}

\vspace{0.3cm}
\begin{block}{Understanding Concentration Accounts}
\begin{itemize}
    \item \textbf{Legitimate use}: Internal bank operations
    \item \textbf{Red flag}: Mixing customer funds
    \item \textbf{Key distinction}: Not for commingling unrelated customers
\end{itemize}
\end{block}
\end{frame}

% ================================
% SLIDE 11: PEP RISK MANAGEMENT
% ================================
\section{PEP Risk Management}
\begin{frame}{PEP Risk Management: Effective Measures}
\begin{block}{Question: Effective Measures for PEP Money Laundering Risks (Select Three)}
\end{block}

\begin{columns}
\column{0.5\textwidth}
\textcolor{wrong}{\textbf{ Wrong:}}
\begin{itemize}
    \item Implement processes to adjust transaction thresholds annually for all clients
\end{itemize}

\column{0.5\textwidth}
\textcolor{correct}{\textbf{Correct:}}
\begin{itemize}
    \item Auto-detect public official occupations
    \item Capture source of funds/wealth at onboarding
    \item Conduct ongoing media monitoring for adverse news
\end{itemize}
\end{columns}

\begin{block}{PEP Risk Management Theory}
\textbf{Prevention + Detection + Ongoing Monitoring}
\end{block}
\end{frame}

% ================================
% SLIDE 12: PEP - THEORY
% ================================
\begin{frame}{PEP Risk Management Theory}
\begin{block}{Three Pillars of PEP Risk Management}
\begin{enumerate}
    \item \textbf{Identification} - Auto-detect public official occupations
    \item \textbf{Assessment} - Capture source of funds and wealth
    \item \textbf{Ongoing Monitoring} - Media screening throughout relationship
\end{enumerate}
\end{block}

\begin{alertblock}{Common Trap}
\textbf{Don't confuse general controls with PEP-specific measures.} \\
Annual threshold adjustments affect all clients - not targeted PEP risk management.
\end{alertblock}
\end{frame}

% ================================
% SLIDE 13: PREPAID CARDS
% ================================
\section{Prepaid Cards}
\begin{frame}{Prepaid Cards: Vulnerability to Financial Crime}
\begin{block}{Question: Why are prepaid cards particularly vulnerable?}
\end{block}

\vspace{0.3cm}
\begin{columns}
\column{0.5\textwidth}
\textcolor{wrong}{\textbf{ Common Misconceptions:}}
\begin{itemize}
    \item "They have built-in controls"
    \item "Limited value makes them low-risk"
\end{itemize}

\column{0.5\textwidth}
\textcolor{correct}{\textbf{Correct Answer:}}
\begin{itemize}
    \item They can be used anonymously, especially when reloadable or unregistered
\end{itemize}
\end{columns}

\begin{block}{Why This Matters}
\begin{itemize}
    \item \textbf{Anonymity} = High ML/TF risk
    \item \textbf{Reloadable} = Can move significant value
    \item \textbf{Cross-border} = Jurisdictional challenges
\end{itemize}
\end{block}
\end{frame}

% ================================
% SLIDE 14: TRADE-BASED MONEY LAUNDERING
% ================================
\section{Trade-Based Money Laundering}
\begin{frame}{Trade-Based Money Laundering (TBML)}
\begin{block}{Question: Which changes reduce TBML exposure? (Select Three)}
\end{block}

\begin{columns}
\column{0.5\textwidth}
\textcolor{wrong}{\textbf{ Wrong:}}
\begin{itemize}
    \item Eliminate letters of credit from jurisdictions with weak AML controls
\end{itemize}

\column{0.5\textwidth}
\textcolor{correct}{\textbf{Correct:}}
\begin{itemize}
    \item Require digitized bills of lading for all high-value shipments
    \item Implement dual-person approval for customer invoice verification
    \item Increase transaction look-back periods in system audits
\end{itemize}
\end{columns}

\begin{block}{TBML Theory}
\textbf{Transparency + Verification + Review}
\end{block}
\end{frame}

% ================================
% SLIDE 15: TBML THEORY
% ================================
\begin{frame}{TBML: Theory and Application}
\begin{block}{Why These Measures Work}
\begin{itemize}
    \item \textbf{Digitized bills of lading}: Prevents document manipulation
    \item \textbf{Dual-person approval}: Segregation of duties
    \item \textbf{Look-back periods}: Pattern detection over time
\end{itemize}
\end{block}

\begin{alertblock}{Avoid This Trap}
\textbf{"Eliminate letters of credit from weak jurisdictions"} \\
This is \textbf{too extreme} - risk-based approach doesn't mean blanket elimination.
\end{alertblock}
\end{frame}

% ================================
% SLIDE 16: SYNTHETIC IDENTITY FRAUD
% ================================
\section{Synthetic Identity Fraud}
\begin{frame}{Synthetic Identity Fraud: Corrective Action}
\begin{block}{The Question}
\small{A new credit product targeting low-income customers shows synthetic identity fraud trends. What is the most appropriate corrective action?}
\end{block}

\vspace{0.3cm}
\begin{columns}
\column{0.5\textwidth}
\textcolor{wrong}{\textbf{ Trap:}}
\begin{itemize}
    \item "Add more verification steps while keeping product active"
    \item "Increase marketing to offset losses"
\end{itemize}

\column{0.5\textwidth}
\textcolor{correct}{\textbf{Correct Answer:}}
\begin{itemize}
    \item Suspend the product while updating risk assessment and onboarding controls
\end{itemize}
\end{columns}

\begin{block}{Why}
\textbf{Risk first, then business} - address root cause before continuing.
\end{block}
\end{frame}

% ================================
% SLIDE 17: GLOBAL AFC STANDARDS - FATF
% ================================
\section{Global AFC Standards}
\begin{frame}{FATF Recommendation 18}
\begin{block}{Question: Deficiencies in Recommendation 18 Implementation}
\small{A multinational institution is criticized for deficiencies in Recommendation 18 implementation. What should auditors have focused on?}
\end{block}

\begin{columns}
\column{0.5\textwidth}
\textcolor{wrong}{\textbf{ Wrong:}}
\begin{itemize}
    \item Training curriculum for non-AML business units
\end{itemize}

\column{0.5\textwidth}
\textcolor{correct}{\textbf{Correct:}}
\begin{itemize}
    \item Whether group-wide AML programs apply to all branches and subsidiaries
\end{itemize}
\end{columns}

\begin{block}{Rec 18: Group-Wide AML Programs}
\textbf{FATF Requirement:} AML programs must apply \textbf{globally} across all branches and subsidiaries.
\end{block}
\end{frame}

% ================================
% SLIDE 18: FSRB vs FATF
% ================================
\begin{frame}{FSRB vs FATF: Understanding the Difference}
\begin{block}{The Question}
\small{A compliance professional compares FATF and FSRB functions while tailoring AML program for a regional subsidiary. Which conclusion best reflects FSRB's role?}
\end{block}

\vspace{0.3cm}
\begin{columns}
\column{0.5\textwidth}
\textcolor{wrong}{\textbf{ Wrong:}}
\begin{itemize}
    \item "FSRBs are separate from FATF"
    \item "FSRBs set independent standards"
\end{itemize}

\column{0.5\textwidth}
\textcolor{correct}{\textbf{Correct Answer:}}
\begin{itemize}
    \item FSRBs help implement and promote FATF standards by evaluating compliance in their regions
\end{itemize}
\end{columns}

\begin{block}{Key Relationship}
\textbf{FSRBs = Regional arms of FATF}, not independent bodies.
\end{block}
\end{frame}

% ================================
% SLIDE 19: FATF EFFECTIVENESS
% ================================
\begin{frame}{FATF Immediate Outcome 5}
\begin{block}{Question: Beneficial Ownership Transparency Weakness}
\small{A country's high volume of legal arrangements relies on anonymous shell entities. National authorities struggle to trace beneficial ownership. Which FATF Immediate Outcome would be rated weak?}
\end{block}

\vspace{0.3cm}
\textcolor{correct}{\textbf{Correct Answer:}}
\begin{center}
\Large{\textbf{Immediate Outcome 5 - Legal persons and arrangements are not misused}}
\end{center}

\begin{block}{Understanding IO5}
\textbf{IO5 measures:} Whether countries prevent misuse of legal entities for money laundering.
\end{block}
\end{frame}

% ================================
% SLIDE 20: BASEL AML INDEX
% ================================
\begin{frame}{Basel AML Index}
\begin{block}{Question: Why Choose the Basel AML Index?}
\small{A national enforcement agency prioritizing policy reform wants to identify vulnerable legal structures contributing to money laundering. Why choose the Basel AML Index?}
\end{block}

\vspace{0.3cm}
\textcolor{correct}{\textbf{Correct Answer:}}
\begin{center}
\Large{It compiles vulnerability indicators across legal, political, and financial dimensions.}
\end{center}

\begin{block}{Understanding the Index}
\begin{itemize}
    \item \textbf{Comprehensive} - Multi-dimensional assessment
    \item \textbf{Comparable} - Cross-country benchmarking
    \item \textbf{Actionable} - Identifies specific vulnerabilities
\end{itemize}
\end{block}
\end{frame}

% ================================
% SLIDE 21: G20 ANTI-CORRUPTION
% ================================
\section{G20 \& International Bodies}
\begin{frame}{G20 Anti-Corruption Working Group}
\begin{block}{Question: G20 ACWG Asset Recovery}
\small{Which G20 ACWG initiative most directly addresses recovering illicitly obtained funds across jurisdictions?}
\end{block}

\begin{columns}
\column{0.5\textwidth}
\textcolor{wrong}{\textbf{ Wrong:}}
\begin{itemize}
    \item Enhancing tax treaties for cross-border information exchange
\end{itemize}

\column{0.5\textwidth}
\textcolor{correct}{\textbf{Correct:}}
\begin{itemize}
    \item Increasing efficiency of asset recovery measures
\end{itemize}
\end{columns}

\begin{block}{Direct vs. Indirect}
\textbf{Asset recovery} is the direct mechanism for recovering funds.
\end{block}
\end{frame}

% ================================
% SLIDE 22: IMF \& WORLD BANK
% ================================
\begin{frame}{IMF/World Bank Expectations}
\begin{block}{Question: Developing Country with Weaknesses}
\small{A developing country received World Bank funding but has ongoing weaknesses in terrorism financing controls and legislative gaps. What aligns with IMF/World Bank expectations?}
\end{block}

\vspace{0.3cm}
\textcolor{correct}{\textbf{Correct Answer:}}
\begin{center}
\Large{Reform its AML/CFT frameworks and align them with international expectations}
\end{center}

\begin{block}{Key Principle}
\textbf{Funding conditional on compliance} - AML/CFT reform is non-negotiable.
\end{block}
\end{frame}

% ================================
% SLIDE 23: IMF IMPROVEMENTS vs CONCERNS
% ================================
\begin{frame}{IMF Assessment: Economic Improvements vs Concerns}
\begin{block}{The Question}
\small{A nation achieved strong GDP growth and reduced inflation after IMF reforms. However, IMF's latest assessment still expresses significant concerns. Which explanation best reflects the IMF's approach?}
\end{block}

\vspace{0.3cm}
\textcolor{correct}{\textbf{Correct Answer:}}
\begin{center}
\Large{Improvements in economic measures can mask financial crime vulnerabilities.}
\end{center}

\begin{block}{Understanding}
\textbf{Strong GDP not equal Strong AML/CFT} - Macroeconomic health doesn't guarantee financial integrity.
\end{block}
\end{frame}

% ================================
% SLIDE 24: IOSCO
% ================================
\begin{frame}{IOSCO: Anti-Financial Crime Objectives}
\begin{block}{Question: IOSCO Objectives Supporting AFC (Select Two)}
\end{block}

\begin{columns}
\column{0.5\textwidth}
\textcolor{wrong}{\textbf{ Wrong:}}
\begin{itemize}
    \item Requiring all securities firms to maintain uniform compliance departments
\end{itemize}

\column{0.5\textwidth}
\textcolor{correct}{\textbf{Correct:}}
\begin{itemize}
    \item Enhancing investor protection through cooperation and enforcement
    \item Promoting financial stability by managing systemic risks and facilitating international information exchange
\end{itemize}
\end{columns}

\begin{block}{Why}
\textbf{Cooperation + Information Exchange + Stability =} Effective AFC in securities.
\end{block}
\end{frame}

% ================================
% SLIDE 25: EGMONT GROUP
% ================================
\section{Egmont Group}
\begin{frame}{Egmont Group Membership Benefits}
\begin{block}{Question: Benefits of Egmont Group Membership (Select Two)}
\end{block}

\begin{columns}
\column{0.5\textwidth}
\textcolor{wrong}{\textbf{ Wrong:}}
\begin{itemize}
    \item Automatic inclusion in FATF mutual evaluations
\end{itemize}

\column{0.5\textwidth}
\textcolor{correct}{\textbf{Correct:}}
\begin{itemize}
    \item Access to secure portal for communication with peer FIUs
    \item Increased capacity through shared typology projects and analyses
\end{itemize}
\end{columns}

\begin{block}{Key Takeaway}
\textbf{Membership not equal Automatic evaluation inclusion} \\
\textbf{Membership = Information sharing capacity}
\end{block}
\end{frame}

% ================================
% SLIDE 26: TECHNOLOGY - AI \& SCREENING
% ================================
\section{Technology for AFC}
\begin{frame}{AI-Based Screening: Addressing Inconsistencies}
\begin{block}{Question: AI Screening System Inconsistencies}
\small{Internal audit reveals AI-based screening system is rejecting transactions inconsistently across business units. Best first step for corrective action?}
\end{block}

\vspace{0.3cm}
\textcolor{correct}{\textbf{Correct Answer:}}
\begin{center}
\Large{Engage the governance committee to assess model performance}
\end{center}

\begin{block}{Why This Approach}
\begin{itemize}
    \item Governance oversight is essential for AI systems
    \item Need centralized assessment of model performance
    \item Inconsistency suggests model governance gap
\end{itemize}
\end{block}
\end{frame}

% ================================
% SLIDE 27: FALSE POSITIVES
% ================================
\begin{frame}{False Positives: Impact and Solution}
\begin{block}{Question: High False Positives in Payment Screening}
\small{A payment screening system returns high false positives. What is likely impact and how should it be addressed?}
\end{block}

\begin{columns}
\column{0.5\textwidth}
\textcolor{wrong}{\textbf{ Wrong:}}
\begin{itemize}
    \item "Reduce list matching criteria to minimize alerts"
\end{itemize}

\column{0.5\textwidth}
\textcolor{correct}{\textbf{Correct:}}
\begin{itemize}
    \item It may cause screening fatigue; tuning scenario logic or data inputs would help
\end{itemize}
\end{columns}

\begin{block}{The Theory}
\textbf{More alerts not equal Better compliance} \\
\textbf{Tuning not equal Reducing controls}
\end{block}
\end{frame}

% ================================
% SLIDE 28: DEVICE INTELLIGENCE
% ================================
\begin{frame}{Device Intelligence Tools}
\begin{block}{Question: Device Intelligence Elements (Select Three)}
\end{block}

\begin{columns}
\column{0.5\textwidth}
\textcolor{wrong}{\textbf{ Wrong:}}
\begin{itemize}
    \item Social media accounts and browsing history
\end{itemize}

\column{0.5\textwidth}
\textcolor{correct}{\textbf{Correct:}}
\begin{itemize}
    \item Browser type and screen resolution
    \item IP address and geolocation
    \item Operating system and installed plugins
\end{itemize}
\end{columns}

\begin{block}{Device Intelligence vs. Social Media}
\textbf{Device intelligence} = Technical device fingerprinting \\
\textbf{Not} = Social media monitoring
\end{block}
\end{frame}

% ================================
% SLIDE 29: ADVERSE MEDIA
% ================================
\begin{frame}{Adverse Media Checks: Why Automate?}
\begin{block}{Question: Why Use Automated Systems for Adverse Media?}
\end{block}

\vspace{0.3cm}
\textcolor{correct}{\textbf{Correct Answer:}}
\begin{center}
\Large{Because they improve efficiency in identifying negative news across multiple sources.}
\end{center}

\begin{block}{The Logic}
\begin{itemize}
    \item Manual review = Limited scope
    \item Automated = Multiple sources simultaneously
    \item Efficiency = More complete coverage
\end{itemize}
\end{block}
\end{frame}

% ================================
% SLIDE 30: AUTOMATED SYSTEMS - MISUSE
% ================================
\begin{frame}{Automated Systems: Incorrect Usage}
\begin{block}{Question: Bank's Automated Systems Not Used Correctly}
\small{A mid-tier bank notices customer transactions flagged as fraudulent. After examination, determined automated systems are not being used correctly. What is most likely cause?}
\end{block}

\vspace{0.3cm}
\textcolor{correct}{\textbf{Correct Answer:}}
\begin{center}
\Large{The bank's automated systems are not being configured to flag transactions that are flagged as fraudulent.}
\end{center}

\begin{block}{Understanding}
\textbf{Systems need proper configuration} - Not just implementation, but ongoing tuning.
\end{block}
\end{frame}

% ================================
% SLIDE 31: REGULATIONS - AI GOVERNANCE
% ================================
\section{AFC Regulations}
\begin{frame}{AI Governance: Regional Variations}
\begin{block}{Question: Regulatory Features of AI Governance (Select Three)}
\end{block}

\begin{columns}
\column{0.5\textwidth}
\textcolor{wrong}{\textbf{ Wrong:}}
\begin{itemize}
    \item Central AI ethics boards with veto power
\end{itemize}

\column{0.5\textwidth}
\textcolor{correct}{\textbf{Correct:}}
\begin{itemize}
    \item Specific AI use-case restrictions
    \item Requirements for human-in-the-loop oversight
    \item Differing standards for transparency and data usage
\end{itemize}
\end{columns}

\begin{block}{Regional Variations}
Different jurisdictions = Different AI governance approaches
\end{block}
\end{frame}

% ================================
% SLIDE 32: MACHINE LEARNING UAE
% ================================
\begin{frame}{Machine Learning in UAE AML}
\begin{block}{Question: Major Advantage of ML Algorithms in UAE AML}
\end{block}

\begin{columns}
\column{0.5\textwidth}
\textcolor{wrong}{\textbf{ Wrong:}}
\begin{itemize}
    \item "They automate product development decisions"
\end{itemize}

\column{0.5\textwidth}
\textcolor{correct}{\textbf{Correct:}}
\begin{itemize}
    \item They improve customer risk assessments by identifying transaction anomalies
\end{itemize}
\end{columns}

\begin{block}{Core Purpose of ML in AML}
\textbf{Detection, not automation of business decisions}
\end{block}
\end{frame}

% ================================
% SLIDE 33: ESG REPORTING
% ================================
\begin{frame}{ESG-Aligned Compliance Requirements}
\begin{block}{Question: ESG Reporting Examples (Select Three)}
\end{block}

\begin{columns}
\column{0.5\textwidth}
\textcolor{wrong}{\textbf{ Wrong:}}
\begin{itemize}
    \item Monitoring account login patterns
    \item Conducting algorithmic risk assessments
\end{itemize}

\column{0.5\textwidth}
\textcolor{correct}{\textbf{Correct:}}
\begin{itemize}
    \item Publishing carbon emissions by facility
    \item Reporting executive salary data across genders
    \item Reporting supply chain labor standards
\end{itemize}
\end{columns}

\begin{block}{Understanding ESG}
\textbf{ESG = Environmental, Social, Governance} \\
Not traditional AML transaction monitoring
\end{block}
\end{frame}

% ================================
% SLIDE 34: EU AML DIRECTIVE
% ================================
\begin{frame}{EU AML Directive Alignment}
\begin{block}{Question: Aligning with Latest EU AML Directive}
\small{European bank wants to align with latest EU AML Directive. What is most effective strategy?}
\end{block}

\begin{columns}
\column{0.5\textwidth}
\textcolor{wrong}{\textbf{ Wrong:}}
\begin{itemize}
    \item Apply FATF guidance uniformly across all non-EU jurisdictions
\end{itemize}

\column{0.5\textwidth}
\textcolor{correct}{\textbf{Correct:}}
\begin{itemize}
    \item Incorporate recommendations from AML regulation updates post-2008 financial crisis
\end{itemize}
\end{columns}

\begin{block}{Key Point}
\textbf{EU-specific} developments matter more than generic FATF guidance.
\end{block}
\end{frame}

% ================================
% SLIDE 35: DATA - EXTERNAL SOURCES
% ================================
\section{Data Collection}
\begin{frame}{External Sources for Sanctions Screening}
\begin{block}{Question: External Sources for Sanctions Screening (Select Two)}
\end{block}

\begin{columns}
\column{0.5\textwidth}
\textcolor{wrong}{\textbf{ Wrong:}}
\begin{itemize}
    \item Watchlists developed from the firm's previous investigations
\end{itemize}

\column{0.5\textwidth}
\textcolor{correct}{\textbf{Correct:}}
\begin{itemize}
    \item FATF blacklists and UN Security Council lists
    \item State Department embargo registries
\end{itemize}
\end{columns}

\begin{block}{External vs Internal}
\textbf{External} = Government/regulatory lists \\
\textbf{Internal} = Proprietary watchlists
\end{block}
\end{frame}

% ================================
% SLIDE 36: DATA LINEAGE
% ================================
\begin{frame}{Data Lineage: Transparency and Auditability}
\begin{block}{Question: Support Transparency through Data Lineage (Select Two)}
\end{block}

\begin{columns}
\column{0.5\textwidth}
\textcolor{wrong}{\textbf{ Wrong:}}
\begin{itemize}
    \item Archiving analyst notes from alert reviews
\end{itemize}

\column{0.5\textwidth}
\textcolor{correct}{\textbf{Correct:}}
\begin{itemize}
    \item Maintaining history of data transformations
    \item Capturing source of each data field
\end{itemize}
\end{columns}

\begin{block}{Data Lineage Definition}
\textbf{Knowing where data comes from and what happens to it.}
\end{block}
\end{frame}

% ================================
% SLIDE 37: BASEL INSTITUTE
% ================================
\section{Basel Institute}
\begin{frame}{Basel Institute: Anti-Corruption Support}
\begin{block}{Question: Basel Institute Services for Anti-Corruption}
\small{NGO developing anti-corruption mechanisms in country with weak legal infrastructure. Which Basel Institute services provide most strategic support?}
\end{block}

\begin{columns}
\column{0.5\textwidth}
\textcolor{wrong}{\textbf{ Wrong:}}
\begin{itemize}
    \item Public finance management technical assistance with standardized compliance templates
\end{itemize}

\column{0.5\textwidth}
\textcolor{correct}{\textbf{Correct:}}
\begin{itemize}
    \item ICAR legal advisory services combined with capacity building programs
\end{itemize}
\end{columns}

\begin{block}{Why}
\textbf{Legal advice + Capacity building} = Sustainable reform
\end{block}
\end{frame}

% ================================
% SLIDE 38: DATA SOURCE CHARACTERISTICS
% ================================
\begin{frame}{Internal vs External Data Sources}
\begin{block}{Question: Internal vs External Source Characteristics (Select Two)}
\end{block}

\vspace{0.3cm}
\textcolor{correct}{\textbf{Correct Answer:}}
\begin{itemize}
    \item Internal sources reflect observed customer behavior or system-held values
    \item External sources include sanction or adverse media lists from agencies
\end{itemize}

\begin{block}{Key Distinction}
\textbf{Internal} = What we know about our customers \\
\textbf{External} = What regulators/agencies tell us
\end{block}
\end{frame}

% ================================
% SLIDE 39: BCBS PEP GUIDANCE
% ================================
\begin{frame}{BCBS Guidance on PEP Escalation}
\begin{block}{Question: Frontline Staff Unsure About PEP Escalation}
\small{During regulatory review, bank discovers frontline staff are unsure when to escalate alerts involving PEPs. Based on BCBS guidance, what is most effective corrective action?}
\end{block}

\vspace{0.3cm}
\textcolor{correct}{\textbf{Correct Answer:}}
\begin{center}
\Large{Reinforce risk-based escalation protocols and provide KYC training}
\end{center}

\begin{block}{BCBS Focus}
\textbf{Training + Clear Protocols} = Effective PEP management
\end{block}
\end{frame}

% ================================
% SLIDE 40: AML HISTORY - EUROPE
% ================================
\begin{frame}{History of AML in Europe}
\begin{block}{Question: European Bank EU AML Directive Alignment}
\small{European bank wants to align with latest EU AML Directive. Given the history of AML in Europe, what is most effective strategy?}
\end{block}

\vspace{0.3cm}
\textcolor{correct}{\textbf{Correct Answer:}}
\begin{center}
\Large{Incorporate recommendations from AML regulation updates post-2008 financial crisis}
\end{center}

\begin{block}{Historical Context}
\textbf{Post-2008} = Major regulatory evolution period
\end{block}
\end{frame}

% ================================
% SLIDE 41: BREACH INVESTIGATION
% ================================
\begin{frame}{Data Breach Investigation}
\begin{block}{Question: Suspicious KYC File Access}
\small{AML compliance officer reviewing system logs after suspected breach shows repeated access to KYC files outside regular business hours by administrative account. Most appropriate next step?}
\end{block}

\vspace{0.3cm}
\textcolor{correct}{\textbf{Correct Answer:}}
\begin{center}
\Large{Investigate the account's activity and suspend its privileges during the review}
\end{center}

\begin{block}{Incident Response}
\textbf{Isolate + Investigate} - Standard security protocol
\end{block}
\end{frame}

% ================================
% SLIDE 42: KEY TAKEAWAYS
% ================================
\section{Key Takeaways}
\begin{frame}{Key Takeaways: Answer Strategy}
\begin{block}{How to Identify Correct Answers}
\begin{enumerate}
    \item \textbf{Specificity} - Prefer specific over general answers
    \item \textbf{Action-Oriented} - Look for concrete actions, not vague concepts
    \item \textbf{Risk-Based} - Choose proportional, risk-based responses
    \item \textbf{Documentation} - Always document, always have audit trail
    \item \textbf{International Standards} - FATF = Default reference point
\end{enumerate}
\end{block}
\end{frame}

% ================================
% SLIDE 43: COMMON TRAPS
% ================================
\begin{frame}{Common Exam Traps to Avoid}
\begin{alertblock}{Most Frequent Mistakes}
\begin{itemize}
    \item \textbf{Over-reliance on automation} - Systems need human oversight
    \item \textbf{Extreme measures} - "Eliminate all" is rarely correct
    \item \textbf{Confusing internal/external} - Know which sources are which
    \item \textbf{General vs specific} - Specific answers usually win
    \item \textbf{Documentation timing} - Record throughout, not just at end
\end{itemize}
\end{alertblock}
\end{frame}

% ================================
% SLIDE 44: THEORY APPLICATION
% ================================
\begin{frame}{Applying Theory to Questions}
\begin{block}{When in Doubt, Ask:}
\begin{itemize}
    \item \textbf{What's the risk?} - Identify actual risk factor
    \item \textbf{What would a regulator expect?} - Regulator perspective
    \item \textbf{What creates an audit trail?} - Documentation focus
    \item \textbf{What's proportional?} - Risk-based approach
    \item \textbf{What does FATF say?} - International standard
\end{itemize}
\end{block}
\end{frame}

% ================================
% SLIDE 45: PRACTICAL TIPS
% ================================
\begin{frame}{Practical Tips for Exam Success}
\begin{columns}
\column{0.5\textwidth}
\begin{block}{Before Answering}
\begin{itemize}
    \item Read carefully
    \item Identify key facts
    \item Note what's missing
    \item Eliminate extremes
\end{itemize}
\end{block}

\column{0.5\textwidth}
\begin{block}{While Answering}
\begin{itemize}
    \item Prefer specific answers
    \item Look for action verbs
    \item Consider auditability
    \item Apply risk-based logic
\end{itemize}
\end{block}
\end{columns}
\end{frame}

% ================================
% SLIDE 46: FINAL SUMMARY
% ================================
\section{Summary}
\begin{frame}{Final Summary}
\begin{block}{Master These Concepts}
\begin{itemize}
    \item \textbf{Transaction Monitoring} - Human + Automated oversight
    \item \textbf{Private Banking} - Discretion and relationship pressure
    \item \textbf{PEP Management} - Identification + Assessment + Monitoring
    \item \textbf{FATF Framework} - IO5, Rec 18, FSRB roles
    \item \textbf{Technology} - AI governance, false positive management
    \item \textbf{Data} - Lineage, internal/external sources
\end{itemize}
\end{block}

\begin{alertblock}{Remember}
\textbf{Think like a compliance professional:}
\textit{"What would I do in this situation?"}
\end{alertblock}
\end{frame}

% ================================
% SLIDE 47: THANK YOU
% ================================
\begin{frame}{Thank You}
\begin{center}
\Huge{\textbf{Good Luck!}}
\vspace{1cm}

\Large{Stay Calm. Think Risk-Based. Apply Theory.}
\vspace{0.5cm}

\normalsize{You've got this!}
\end{center}
\end{frame}

\end{document}